	\documentclass[14pt]{extreport}
\usepackage{extsizes}
	\usepackage[french]{babel}
	\usepackage[utf8]{inputenc}  
	\usepackage[T1]{fontenc}
	\usepackage{amssymb}
	\usepackage[mathscr]{euscript}
	\usepackage{stmaryrd}
	\usepackage{amsmath}
	\usepackage{tikz}
	\usepackage[all,cmtip]{xy}
	\usepackage{amsthm}
	\usepackage{varioref}
	\usepackage[ margin=0.6in]{geometry}
	\geometry{a4paper}
	\usepackage{lmodern}
	\usepackage{hyperref}
	\usepackage{array}
	\usepackage{float}
	\usepackage{easytable}
	 \usepackage{fancyhdr}\usepackage{longtable}
	 \usetikzlibrary{shapes.misc}
	 \newcommand\ang[1]{$#1{}^o$}
\newlength{\taillecellule}
\setlength{\taillecellule}{2cm}
\newcolumntype{C}{@{}>{\centering\arraybackslash}p{\taillecellule}@{}}
 

	\pagestyle{fancy}
	\theoremstyle{plain}
	\fancyfoot[C]{\empty} 
	\fancyhead[L]{Interrogation chapitre 8}
	\fancyhead[R]{3 juin 2024}
	
	
	\title{Interrogation chapitre 7}
	\date{}
	
	
	\begin{document}
 \subsection*{Exercice 1 : 4 points}
Calculez, en détaillant les étapes sur votre copie : 

\begin{enumerate}
\item $2 - (-3) + 5 \times  2$
\item $ -2 + 3 \times 4$
\item $((5 - 4) \times 2 + 3 )\times 7 - 4$
\item $-2 - (3 - 1)$.
\end{enumerate}

\subsection*{Exercice 2 : 4 points}
Évaluez les expressions suivantes en la valeur indiquée : 
\begin{enumerate}
\item L'expression $x+5 - 1 $ en $x=9$
\item L'expression $(x+5)\times 3$ en $x=-1$.
\item L'expression $3x-2$ en $x=7$.
\item L'expression $3(x-1)+4x$ en $x=+2$.
\end{enumerate}

\subsection*{Exercice 3 (1+2+1+1+1+1= 7 points)}

On considère la figure ci-dessous, où $D$ est le milieu de $[AB]$, $E$ le milieu de $[AC]$ et où $G'$ est le symétrique de $G$ par rapport à $D$.
\begin{enumerate}
\item Montrer que $D$ est le milieu de $[GG']$. 
\item En déduire la nature de $GBG'C$. 
\item En déduire que $(GB)//(CG')$. 
\item À l'aide du théorème des milieux, montrer que $G$ est le milieu de $[AG']$. 
\item En déduire que $AG'= 4\times GD$
\item Montrer que $AD= 3\times GD$. (On a montré que les médianes d'un triangle se coupent à un tiers de leurs longueurs).
\end{enumerate}

\includegraphics[scale=.5]{Exo3}

\subsection*{Exercice 4 (2 + 1 + 2 = 5 points)} % 7,5 pts
On considère un parallélogramme $ABCD$, et on note $O$ l'intersection des diagonales $[AC]$ et $[BD]$. On place sur la diagonale $[AC]$ les points $I$ et $J$ la coupant en trois parties égales : $AI=IJ=JC$.  On suppose $AB = 8$ cm, $AC = 6$ cm et $BC = 4$ cm.

\includegraphics[scale=.45]{Exo4}
\begin{enumerate} 
\item Calculer les longueurs $OA$ et $OC$, en détaillant votre raisonnement.
\item En déduire les longueurs $OI$ et $OJ$.
\item En déduire que $IBJD$ est un parallélogramme.  
\item [Bonus 4.] Montrer que les triangles $AID$ et $BJC$ sont égaux (On pourra utiliser que deux triangles qui ont les mêmes longueurs sont égaux).
\end{enumerate}
\end{document}

